\chapter{方法}
\label{ch:methodology-zh}
\label{ch:architecture-zh}

本章呈現 RA-GARK 之完整方法。第~\ref{sec:approach-principle-zh}~節說明使 RA-GARK 有別於第~\ref{ch:relatedwork}~章所回顧之深度融合基線的核心設計原則。第~\ref{sec:approach-overview-zh}~節以泳道圖（swim-lane）呈現推論時的整體資料流。第~\ref{sec:approach-notation-zh}~節固定符號。第~\ref{sec:approach-problem-zh}~節形式化推薦任務與訓練目標。第~\ref{sec:local-zh}~至~\ref{sec:cl-zh}~節依模組於前向傳遞中的順序逐一詳述各模組；第~\ref{sec:method-train-zh}~節說明訓練程序與計算複雜度。

\section{設計原則：KG 作為可被閘控的旁路通道}
\label{sec:approach-principle-zh}

第~\ref{ch:intro}~章的經驗觀察是：在本研究稀疏的評論衍生知識圖譜（KG）上，四種具代表性的 KG-aware 方法（KGAT、KGCL、MCCLK、KGRec）皆被純 LightGCN 一致地超越。這些基線所共有的結構性失因在於：KG 被當作評分管線的\emph{必要組件}串接進來——KG 實體嵌入直接參與訊息傳遞，且無任何架構性的開關能在 KG 被證實不可靠時將其脫接。因此 KG 雜訊與協同表示共享同一條梯度流。

RA-GARK 圍繞一個直接面對此失效模式的核心設計原則：

\begin{quote}
\emph{KG 不應為評分管線的必要組件，而應為一條旁路通道；其對最終分數的貢獻由一個學習得到的閘決定，閘在初始化時偏置至一個安全預設（純協同過濾），僅在訓練訊號顯示 KG 通道能降低損失時才被打開。}
\end{quote}

此原則衍生三項架構層級的後果，並共同定義 RA-GARK：

\begin{enumerate}
  \item \textbf{雙條平行視角。} \emph{區域視角}透過在使用者—物品互動圖上運行 LightGCN 產生純協同表示 $(u_{\mathrm{loc}}, i_{\mathrm{loc}})$，過程中完全不涉及 KG。\emph{全域視角}則由獨立的使用者嵌入表與物品端的潛在面向槽產生 KG 表示 $(u_{\mathrm{glo}}, i_{\mathrm{glo}})$，過程中亦不涉及協同訊號。直至融合階段之前，兩條視角在結構上完全分離。
  \item \textbf{後期、閘控的融合。} 兩條視角僅在最終評分階段被融合閘合併；該閘針對每組（使用者, 物品）配對輸出一個混合權重 $\alpha \in (0, 1)$。閘以偏置初始化使 $\alpha^{(0)} \approx 0.993$，模型訓練起點實質等同於純 LightGCN，KG 通道僅在梯度將 $\alpha$ 推低時逐步開啟。
  \item \textbf{合理化感知的物品端聚合。} 每件物品的 KG 語意儲存於 $A=4$ 個\emph{潛在}面向槽中，槽位內容由 IDF 加權商品—面向矩陣之截斷 SVD 產生；softmax aspect-saliency 注意力依（使用者, 物品）配對選取代表該物品的潛在槽。
\end{enumerate}

每項後果均於第~\ref{ch:methodology}~章中數學化。本章其餘部分將符號與問題陳述定位至足以使第~\ref{ch:methodology}~章自洽所需之細節層級。

\section{架構概覽}
\label{sec:approach-overview-zh}

圖~\ref{fig:architecture-zh} 勾勒了推論時的資料流。模型被布置為對應區域與全域視角的兩條水平泳道，於右側由融合階段匯合。區域視角（上方泳道）對使用者—物品互動圖套用 LightGCN，得到 $(u_{\mathrm{loc}}, i_{\mathrm{loc}})$。全域視角（下方泳道）直接以查表得到使用者的全域嵌入 $u_{\mathrm{glo}}$，並對物品的 $A=4$ 個潛在面向槽施加 softmax 合理化遮罩以得 $i_{\mathrm{glo}}$；面向槽本身為 SVD 初始化（圖中以虛線標註）。兩條視角由獨立的使用者側與物品側融合閘混合，每個閘的偏置初始化為 $+5$；最終分數為內積 $\langle u_{\mathrm{final}}, i_{\mathrm{final}} \rangle$。為求清晰，跨視角對比學習路徑於圖中省略，詳見第~\ref{sec:cl-zh}~節。

\begin{figure}[ht]
\centering
\includegraphics[width=\textwidth]{img/architecture.png}
\caption{RA-GARK 架構。}
\label{fig:architecture-zh}
\end{figure}

兩項結構性質值得在形式化呈現之前先點明。其一，$u_{\mathrm{glo}}$ 並不經過合理化遮罩模組，而是由直接查表得到，並逕送至使用者側融合閘；合理化遮罩僅作用於物品側，符合「合理化回答的是『此物品的哪個面向對此使用者重要』」這一語意。其二，KG-SVD 步驟為初始化而非前向運算：SVD 於訓練前計算一次以設定面向槽參數的初值，其後槽位即如其他參數般以梯度下降更新。

\section{符號}
\label{sec:approach-notation-zh}

本論文後續使用之符號列於表~\ref{tab:notation-zh}。

\begin{table}[ht]
\centering
\caption{本章與後續章節所使用之符號。}
\label{tab:notation-zh}
\begin{tabular}{ll}
\toprule
符號 & 意義 \\
\midrule
$\mathcal{U}, \mathcal{I}$           & 使用者集合與物品集合 \\
$\mathcal{R} \subseteq \mathcal{U} \times \mathcal{I}$ & 已觀察到的使用者—物品互動 \\
$\mathcal{G}_{\mathrm{KG}}$          & 物品—面向知識圖譜（評論衍生） \\
$\mathcal{A}$                        & $\mathcal{G}_{\mathrm{KG}}$ 中出現之 unique 面向標籤集合 \\
$A$                                  & 每件物品的潛在面向槽數（$A=4$） \\
$d$                                  & 嵌入維度（除另行指定，$d=64$） \\
$K$                                  & LightGCN 傳播層數（$K=2$） \\
$u_{\mathrm{loc}}, i_{\mathrm{loc}} \in \mathbb{R}^d$ & 區域視角（LightGCN）嵌入 \\
$u_{\mathrm{glo}} \in \mathbb{R}^d$   & 全域視角使用者嵌入（直接查表） \\
$\mathbf{a}_i \in \mathbb{R}^{A \times d}$ & 物品 $i$ 之潛在面向槽（SVD 初始化） \\
$i_{\mathrm{glo}} \in \mathbb{R}^d$    & 全域視角物品嵌入（合理化遮罩後） \\
$\alpha_u, \alpha_i \in (0,1)$        & 使用者側與物品側融合閘輸出 \\
$\hat{y}(u, i) \in \mathbb{R}$         & 對配對 $(u, i)$ 之預測偏好分數 \\
\bottomrule
\end{tabular}
\end{table}

\section{問題定式}
\label{sec:approach-problem-zh}

\paragraph{任務。} 本文處理隱式回饋設定下的 top-$K$ 推薦。給定 $\mathcal{U}$、$\mathcal{I}$、已觀察到的互動 $\mathcal{R}$，以及物品—面向 KG $\mathcal{G}_{\mathrm{KG}}$，目標是為每位使用者 $u \in \mathcal{U}$ 從 $\mathcal{I} \setminus \{i : (u, i) \in \mathcal{R}\}$ 中產生一個排序清單，並以排序指標（NDCG@$K$、Recall@$K$、Hit@$K$）對保留之測試互動集進行評估。

\paragraph{評分函數。} RA-GARK 對每組 $(u, i)$ 配對指派內積分數
\begin{equation}
\label{eq:score-zh}
\hat{y}(u, i) = \langle u_{\mathrm{final}},\ i_{\mathrm{final}} \rangle,
\end{equation}
其中 $u_{\mathrm{final}}$ 與 $i_{\mathrm{final}}$ 由區域與全域表示的閘控融合產生：
\begin{equation}
\label{eq:fusion-zh}
u_{\mathrm{final}} = \alpha_u\, u_{\mathrm{loc}} + (1 - \alpha_u)\, u_{\mathrm{glo}}, \qquad
i_{\mathrm{final}} = \alpha_i\, i_{\mathrm{loc}} + (1 - \alpha_i)\, i_{\mathrm{glo}}.
\end{equation}
閘輸出 $\alpha_u, \alpha_i$ 分別由獨立的 MLP 計算（第~\ref{sec:gate}~節）。在本論文採用的偏置初始化下，$\alpha_u^{(0)}, \alpha_i^{(0)} \approx 0.993$，因此 $\hat{y}^{(0)}(u, i) \approx \langle u_{\mathrm{loc}}, i_{\mathrm{loc}} \rangle$，即訓練起點之分數實質為純 LightGCN 之分數，直至閘藉學習而調整。

\paragraph{訓練目標。} 我們以 Bayesian Personalized Ranking（BPR）~\cite{rendle2009bpr} 損失為主體，並輔以兩條跨視角對比正則化：
\begin{equation}
\label{eq:total-loss-zh}
\mathcal{L} = \mathcal{L}_{\mathrm{BPR}} + \lambda_{\mathrm{CL}} \bigl( \mathcal{L}_{\mathrm{aCL}} + \mathcal{L}_{\mathrm{uCL}} \bigr),
\end{equation}
其中對每組已觀察配對 $(u, i^+) \in \mathcal{R}$ 與一個取樣的負例物品 $i^- \notin \mathcal{R}_u$，
\begin{equation}
\label{eq:bpr-zh}
\mathcal{L}_{\mathrm{BPR}} = - \log \sigma\bigl( \hat{y}(u, i^+) - \hat{y}(u, i^-) \bigr),
\end{equation}
而 $\mathcal{L}_{\mathrm{aCL}}, \mathcal{L}_{\mathrm{uCL}}$ 分別為 $(i_{\mathrm{loc}}, i_{\mathrm{glo}})$ 與 $(u_{\mathrm{loc}}, u_{\mathrm{glo}})$ 之間的 InfoNCE 式對齊損失；全域側施加 stop-gradient 以避免對比梯度擾動 SVD 初始化之面向槽（第~\ref{sec:cl}~節）。對比權重 $\lambda_{\mathrm{CL}}$ 固定為 $0.005$，使對比的角色為輔助而非主要。

\section{區域視角：LightGCN}
\label{sec:local-zh}

% TODO: 區域視角內文
% - K=2 LightGCN 傳播 + 各層平均
% - 不涉及 KG；作為協同訊號純淨化錨點
% - 引用 he2020lightgcn

\section{全域視角——KG-SVD 初始化}
\label{sec:kgsvd-zh}

% TODO: KG-SVD 初始化內文
% - 物品—面向共現矩陣 M
% - IDF 加權動機
% - 截斷 SVD，k = A * d
% - reshape 至 [N_i, A, d]
% - xavier 等量縮放
% - 消融：移除 SVD init → −5.3% NDCG@20

\begin{figure}[ht]
  \centering
  % \includegraphics[width=0.9\textwidth]{img/kgsvd_init.png}  % TODO: 插入圖 (b)
  \fbox{\parbox[c][4cm][c]{0.9\textwidth}{\centering\itshape 佔位：KG-SVD 初始化流程\\(物品—面向矩陣 $\to$ IDF 加權 $\to$ 截斷 SVD $\to$ reshape 至 $[N_i, A, d]$)。}}
  \caption{KG-SVD 初始化流程。}
  \label{fig:kgsvd-zh}
\end{figure}

\section{全域視角——Softmax 合理化遮罩}
\label{sec:softmax-zh}

% TODO: softmax 合理化遮罩內文
% - MLP([u_glo; a_{i,k}]) 每槽位的 logit
% - 對 A 個槽位的 softmax，溫度 τ=0.5
% - 加權求和形成 i_glo
% - 合理化 = 此物品的哪個潛在槽位代表給定的使用者—物品配對
% - softmax vs sigmoid 的設計理由（競爭 vs 獨立）
% - 消融：sigmoid → −7.0% NDCG@20

\begin{figure}[ht]
  \centering
  % \includegraphics[width=0.9\textwidth]{img/softmax_attention.png}  % TODO: 插入圖 (c)
  \fbox{\parbox[c][4cm][c]{0.9\textwidth}{\centering\itshape 佔位：Softmax aspect-saliency 注意力前向流程\\($u_{\mathrm{glo}}$, $\{a_{i,k}\}_{k=1}^A$ $\to$ MLP logits $\to$ softmax 權重 $\to$ 加權求和 $\to$ $i_{\mathrm{glo}}$)。}}
  \caption{Softmax 合理化遮罩。}
  \label{fig:softmax-zh}
\end{figure}

\section{區域偏置融合閘}
\label{sec:gate-zh}

% TODO: 融合閘內文
% - 獨立之 user_fusion_gate 與 item_fusion_gate
% - MLP([x_loc; x_glo]) → tanh → linear → sigmoid → α
% - 最後線性層偏置初始化 +5; α₀ ≈ σ(5) ≈ 0.993
% - 優雅退化：模型起點 ≈ 純 LightGCN
% - 消融：bias=0 → −5.3% NDCG@20

\begin{figure}[ht]
  \centering
  % \includegraphics[width=0.9\textwidth]{img/fusion_gate.png}  % TODO: 插入圖 (d)
  \fbox{\parbox[c][4cm][c]{0.9\textwidth}{\centering\itshape 佔位：融合閘訊號流\\($x_{\mathrm{loc}}$, $x_{\mathrm{glo}}$ $\to$ 偏置初始化最後線性層 ($b=+5$) 之 MLP $\to$ sigmoid $\to$ $\alpha$ $\to$ $x_{\mathrm{final}} = \alpha\, x_{\mathrm{loc}} + (1-\alpha)\, x_{\mathrm{glo}}$)。使用者與物品各一個獨立閘。}}
  \caption{區域偏置融合閘。}
  \label{fig:gate-zh}
\end{figure}

\section{跨視角對比正則化}
\label{sec:cl-zh}

% TODO: 對比學習內文
% - L_aCL between (i_loc, i_glo); L_uCL between (u_loc, u_glo)
% - InfoNCE 形式 + cl_projector 投影頭
% - 全域側 stop-gradient
% - 小固定權重 λ_CL = 0.005（輔助角色）

\section{訓練與計算複雜度}
\label{sec:method-train-zh}

% TODO: 訓練與複雜度內文
% - Adam, lr, batch size, early stopping
% - 各模組複雜度
% - 全項目排序之推論成本
